\section{Case study: three global models over Germany}
\label{sec:case}

The case study measures how quickly skill decays, what simple post-processing adds, and whether the four propositions of Section~\ref{sec:foundations} hold on real forecasts. It uses only public data, and every number below is written by the analysis code (Appendix~\ref{app:repro}).

\subsection{Data and design}

\para{Forecasts.} The 00\,UTC runs of 1 May to 31 July 2025 (\valNumRuns{} runs per model) of ECMWF IFS and ECMWF AIFS Single from ECMWF open data~\cite{ecmwf2025opendata}, and of NOAA GFS from the NOAA archive~\cite{noaa2025gfs}, all on 0.25$^\circ$ grids, at lead times of 0 to 240\,h in 12-hour steps. For each run and lead time only the 10\,m wind components and 2\,m temperature were downloaded, with byte-range requests on the GRIB2 indexes, and interpolated bilinearly to the stations; every download is logged with its SHA-256 hash.

\para{Observations.} Hourly reports of the NOAA Integrated Surface Database~\cite{smith2011isd} from German stations below 600\,m. The report closest to the valid time within 30 minutes is used; stations are kept if both variables are available at 90\,\% or more of the valid times, which leaves \valNumStations{} stations, mostly airports, spread over the whole country.

\para{Forecasts compared.} The three raw models; their multi-model mean (MMM), the plain average; the same four after a running bias correction (BC) that subtracts, for each model, station and lead time, the mean error of the previous 30 runs whose observations were already available when the run started; and two references: persistence, the latest observation at the same time of day available at the start of the run, and the 2020--2024 station climatology for the same month and hour, which is independent of the test data.

\para{Verification.} Scores are computed for the \valNumEvalRuns{} runs from 1 June (May only spins up the bias correction), giving \valNumEvalPairs{} model-observation pairs per variable. Intervals are 95\,\% bootstrap intervals that resample whole runs, which keeps the correlation between stations~\cite{kunsch1989bootstrap}; differences are tested with the Diebold-Mariano test on daily mean squared errors~\cite{diebold1995comparing}. Valid times alternate between 12\,UTC (half-day lead times) and 00\,UTC (whole days); the figures show the 12\,UTC valid times, when solar output and wind ramps matter most, and the data files contain both.

\subsection{Skill decays within days, and post-processing buys days}

% Figure 5: skill by lead time. Data: data/processed/fig05_skill_by_lead_{wind,t2m}.csv
% Script: code/04_skill_by_lead_time.py. Only 12 UTC valid times are drawn.
\begin{figure*}[t]
\centering
\readcsv{\datadir/fig05_skill_by_lead_wind.csv}\skwind
\readcsv{\datadir/fig05_skill_by_lead_t2m.csv}\sktemp
\pgfplotsset{
  noon/.style={only if={valid_hour_utc}{12}},
  ifs/.style={ewfC1,line width=1.1pt,mark=*,mark size=1.3pt},
  aifs/.style={ewfC2,line width=1.1pt,mark=*,mark size=1.3pt},
  gfs/.style={ewfC3,line width=1.1pt,mark=*,mark size=1.3pt},
  mmmraw/.style={ewfInkMuted,line width=1.1pt,mark=*,mark size=1.3pt},
  mmmbc/.style={ewfInk,line width=1.6pt,mark=*,mark size=1.6pt},
}
\begin{tikzpicture}
\begin{groupplot}[group style={group size=2 by 2, horizontal sep=1.6cm, vertical sep=1.55cm},
  paper, width=0.47\textwidth, height=5.1cm, xmin=0, xmax=10, xtick={0,1,...,10}]
% ---------------- (a) wind RMSE
\nextgroupplot[title={(a) 10\,m wind speed: RMSE}, ylabel={RMSE (m/s)}, ymin=1.0, ymax=3.0,
  legend to name=skilllegend, legend columns=4,
  legend style={font=\scriptsize,/tikz/every even column/.append style={column sep=0.45cm}}]
\addplot[name path=hi,draw=none,forget plot,noon] table[x=lead_days,y=rmse_hi_mmm_bc]{\skwind};
\addplot[name path=lo,draw=none,forget plot,noon] table[x=lead_days,y=rmse_lo_mmm_bc]{\skwind};
\addplot[band=ewfInk,forget plot] fill between[of=hi and lo];
\addplot[ifs,noon] table[x=lead_days,y=rmse_ifs]{\skwind}; \addlegendentry{ECMWF IFS}
\addplot[aifs,noon] table[x=lead_days,y=rmse_aifs]{\skwind}; \addlegendentry{ECMWF AIFS}
\addplot[gfs,noon] table[x=lead_days,y=rmse_gfs]{\skwind}; \addlegendentry{NOAA GFS}
\addplot[mmmraw,noon] table[x=lead_days,y=rmse_mmm]{\skwind}; \addlegendentry{multi-model mean (MMM)}
\addplot[mmmbc,noon] table[x=lead_days,y=rmse_mmm_bc]{\skwind}; \addlegendentry{MMM, bias-corrected}
\addplot[reference,noon] table[x=lead_days,y=rmse_climatology]{\skwind}; \addlegendentry{climatology 2020--2024}
\addplot[reference,dotted,line width=0.9pt,noon] table[x=lead_days,y=rmse_persistence]{\skwind}; \addlegendentry{persistence}
% ---------------- (b) wind skill
\nextgroupplot[title={(b) 10\,m wind speed: skill against climatology}, ylabel={MSE skill score}, ymin=-0.8, ymax=0.8]
\addplot[reference,domain=0:10,samples=2,forget plot] {0};
\addplot[ifs,noon] table[x=lead_days,y=skill_ifs]{\skwind};
\addplot[aifs,noon] table[x=lead_days,y=skill_aifs]{\skwind};
\addplot[gfs,noon] table[x=lead_days,y=skill_gfs]{\skwind};
\addplot[mmmraw,noon] table[x=lead_days,y=skill_mmm]{\skwind};
\addplot[mmmbc,noon] table[x=lead_days,y=skill_mmm_bc]{\skwind};
\node[font=\tiny,text=ewfInkSecondary,anchor=south west] at (axis cs:0.1,0.0) {better than climatology};
% ---------------- (c) temperature RMSE
\nextgroupplot[title={(c) 2\,m temperature: RMSE}, ylabel={RMSE (K)}, ymin=0.5, ymax=6.5, xlabel={Lead time (days)}]
\addplot[name path=hi,draw=none,forget plot,noon] table[x=lead_days,y=rmse_hi_mmm_bc]{\sktemp};
\addplot[name path=lo,draw=none,forget plot,noon] table[x=lead_days,y=rmse_lo_mmm_bc]{\sktemp};
\addplot[band=ewfInk,forget plot] fill between[of=hi and lo];
\addplot[ifs,noon] table[x=lead_days,y=rmse_ifs]{\sktemp};
\addplot[aifs,noon] table[x=lead_days,y=rmse_aifs]{\sktemp};
\addplot[gfs,noon] table[x=lead_days,y=rmse_gfs]{\sktemp};
\addplot[mmmbc,noon] table[x=lead_days,y=rmse_mmm_bc]{\sktemp};
\addplot[reference,noon] table[x=lead_days,y=rmse_climatology]{\sktemp};
\addplot[reference,dotted,line width=0.9pt,noon] table[x=lead_days,y=rmse_persistence]{\sktemp};
% ---------------- (d) temperature skill
\nextgroupplot[title={(d) 2\,m temperature: skill against climatology}, ylabel={MSE skill score}, ymin=-0.4, ymax=1.0, xlabel={Lead time (days)}]
\addplot[reference,domain=0:10,samples=2,forget plot] {0};
\addplot[ifs,noon] table[x=lead_days,y=skill_ifs]{\sktemp};
\addplot[aifs,noon] table[x=lead_days,y=skill_aifs]{\sktemp};
\addplot[gfs,noon] table[x=lead_days,y=skill_gfs]{\sktemp};
\addplot[mmmraw,noon] table[x=lead_days,y=skill_mmm]{\sktemp};
\addplot[mmmbc,noon] table[x=lead_days,y=skill_mmm_bc]{\sktemp};
\end{groupplot}
\end{tikzpicture}\\[3pt]
\ref{skilllegend}
\caption{\textbf{Skill decays within days; averaging and bias correction buy back about three.} Verification at 12\,UTC valid times for \valNumEvalRuns{} runs (June--July 2025) at \valNumStations{} German stations. MMM: multi-model mean of IFS, AIFS and GFS; BC: running 30-day bias correction. The grey band is the 95\,\% bootstrap interval of the bias-corrected MMM. The skill score is $1-\text{MSE}/\text{MSE}_{\text{climatology}}$ with the 2020--2024 station climatology; below zero the forecast is worse than climatology.}
\label{fig:skill}
\end{figure*}


Figure~\ref{fig:skill} and Table~\ref{tab:da} summarise the results. At day-ahead lead times the raw IFS and AIFS forecasts of 10\,m wind are statistically indistinguishable (RMSE \valWindIfsRmseDayAhead{} and \valWindAifsRmseDayAhead\,m/s, Diebold-Mariano $p=\valWindDmAifsvsIfsP$), GFS is worse (\valWindGfsRmseDayAhead\,m/s, $p\valWindDmGfsvsIfsP$), and all three beat persistence (\valWindPersistenceRmseDayAhead\,m/s) and climatology (\valWindClimatologyRmseDayAhead\,m/s). For 2\,m temperature AIFS is slightly better than IFS (\valTempAifsRmseDayAhead{} vs.\ \valTempIfsRmseDayAhead\,K, $p=\valTempDmAifsvsIfsP$) and GFS clearly worse.

The raw wind forecasts carry a large systematic error with a day-night cycle: at 12\,UTC on the second day the models underestimate the observed wind by \valWindIfsBiasThirtySix{} (IFS) and \valWindAifsBiasThirtySix{}\,m/s (AIFS), at 00\,UTC they overestimate it by about \valWindIfsBiasTwentyFour\,m/s, a typical signature of near-surface stability and of the difference between a 0.25$^\circ$ grid box and an airport anemometer. The AIFS underestimate grows with lead time, to \valWindAifsBiasNineHalfDay\,m/s at 9.5 days, consistent with the smoothing of deterministic machine-learning forecasts. The 30-day bias correction removes most of this error: it lowers the day-ahead wind RMSE of the multi-model mean by \valWindBcGainPct\,\%, from \valWindMmmRmseDayAhead{} to \valWindMmmBcRmseDayAhead\,m/s ($p\valWindDmMmmBcvsMmmP$), and the temperature RMSE by \valTempBcGainPct\,\%.

Skill against climatology (Figure~\ref{fig:skill}b) defines the useful horizon. For daytime 10\,m wind it reaches zero after \valWindGfsSkillZeroDays{} days for GFS, \valWindIfsSkillZeroDays{} for IFS and \valWindAifsSkillZeroDays{} for AIFS; averaging extends it to \valWindMmmSkillZeroDays{} days and bias correction of the average to \valWindMmmBcSkillZeroDays{} days. Temperature keeps positive skill for longer, beyond 9.5 days for AIFS and the multi-model mean. Model averaging and a simple bias correction together thus extend the useful daytime wind forecast by about three days, without any change to the weather models.

% Table 7: day-ahead verification. The tabular is generated by code/04_skill_by_lead_time.py
% into tab07-day-ahead-summary.tex; data: data/processed/tab07_day_ahead_summary.csv
\begin{table*}[t]
\centering\small
\caption{\textbf{Day-ahead verification against \valNumStations{} German stations} (lead times 24, 36 and 48\,h; \valNumEvalRuns{} runs from 1 June to 31 July 2025). Bias is forecast minus observation. BC: running 30-day bias correction using past data only. Intervals are 95\,\% bootstrap intervals over runs. Data: \texttt{data/processed/tab07\_day\_ahead\_summary.csv}.}
\label{tab:da}
\setlength{\tabcolsep}{5pt}
% Generated by code/04_skill_by_lead_time.py -- do not edit by hand.
\begin{tabular}{@{}lrrrrrr@{}}
\toprule
& \multicolumn{3}{c}{10\,m wind speed (m/s)} & \multicolumn{3}{c}{2\,m temperature (K)} \\
\cmidrule(lr){2-4}\cmidrule(l){5-7}
Forecast & Bias & MAE & RMSE [95\,\% CI] & Bias & MAE & RMSE [95\,\% CI] \\
\midrule
ECMWF IFS & $0.02$ & 1.06 & 1.37 [1.32, 1.42] & $0.08$ & 1.19 & 1.56 [1.49, 1.64] \\
ECMWF AIFS & $-0.07$ & 1.06 & 1.36 [1.31, 1.41] & $0.32$ & 1.15 & 1.51 [1.43, 1.59] \\
NOAA GFS & $0.24$ & 1.14 & 1.46 [1.40, 1.52] & $-0.67$ & 1.62 & 2.05 [1.98, 2.13] \\
Multi-model mean & $0.06$ & 1.02 & 1.30 [1.25, 1.35] & $-0.09$ & 1.12 & 1.46 [1.40, 1.51] \\
\addlinespace
ECMWF IFS (BC) & $-0.03$ & 0.91 & 1.20 [1.15, 1.24] & $-0.01$ & 1.15 & 1.51 [1.45, 1.58] \\
ECMWF AIFS (BC) & $-0.05$ & 0.87 & 1.14 [1.10, 1.19] & $-0.04$ & 1.04 & 1.35 [1.30, 1.40] \\
NOAA GFS (BC) & $-0.03$ & 0.98 & 1.28 [1.22, 1.34] & $0.17$ & 1.34 & 1.75 [1.68, 1.83] \\
\textbf{Multi-model mean (BC)} & $-0.04$ & 0.86 & 1.12 [1.08, 1.16] & $0.04$ & 1.03 & 1.35 [1.30, 1.41] \\
\addlinespace
Persistence & $-0.03$ & 1.62 & 2.16 & $0.12$ & 3.17 & 4.16 \\
Climatology 2020--2024 & $-0.24$ & 1.25 & 1.66 & $-0.37$ & 2.67 & 3.47 \\
\bottomrule
\end{tabular}

\end{table*}


\para{Model averaging (Proposition~\ref{prop:combine}).} The day-ahead wind errors of IFS and AIFS have a correlation of \valWindIfsAifsErrCorrDayAhead{} and an RMSE ratio of \valWindIfsAifsRatioDayAhead, so $k^2+2\rho k$ lies just below 3 and averaging helps, but only a little: the three-model mean is \valWindMmmGainPct\,\% better than the best single model. The error correlation falls with lead time, from \valWindRawErrCorrTwelve{} at 12\,h to \valWindRawErrCorrTenDay{} at 10 days, so the benefit of averaging grows with lead time, as Figure~\ref{fig:skill} shows. After bias correction AIFS pulls ahead of IFS (RMSE ratio up to \valWindBcMaxRatio), and at \valWindBcCombFailLeads{} of \valNumLeadTimes{} wind lead times the condition fails: there the IFS-AIFS average is worse than AIFS alone, exactly as the proposition predicts.

\subsection{Errors that average out, and errors that do not}

% Figure 6: spatial aggregation. Data: data/processed/fig06a_*.csv, fig06b_aggregation_rmse_vs_n.csv
% Script: code/05_spatial_aggregation.py
\begin{figure*}[t]
\centering
\readcsv{\datadir/fig06a_error_correlation_pairs.csv}\aggpairs
\readcsv{\datadir/fig06a_error_correlation_binned.csv}\aggbins
\readcsv{\datadir/fig06b_aggregation_rmse_vs_n.csv}\aggn
\begin{tikzpicture}
\begin{groupplot}[group style={group size=2 by 1, horizontal sep=1.6cm},
  paper, width=0.47\textwidth, height=5.2cm]
% ---------------- (a)
\nextgroupplot[title={(a) Error correlation between station pairs}, xlabel={Distance (km)},
  ylabel={Correlation of day-ahead errors}, xmin=0, xmax=850, ymin=-0.3, ymax=0.7,
  legend style={at={(0.98,0.98)},anchor=north east}]
\addplot[reference,domain=0:850,samples=2,forget plot] {0};
\addplot[only marks,mark=*,mark size=0.9pt,ewfC1,fill opacity=0.35,draw opacity=0] table[x=distance_km,y=error_correlation]{\aggpairs}; \addlegendentry{station pair}
\addplot[ewfInk,line width=1.4pt,mark=*,mark size=1.8pt] table[x=distance_mid_km,y=mean_correlation]{\aggbins}; \addlegendentry{mean per 100\,km}
% ---------------- (b)
\nextgroupplot[title={(b) Error of a portfolio of $N$ stations}, xlabel={Number of stations $N$},
  ylabel={RMSE relative to one station}, xmin=1, xmax=62, ymin=0, ymax=1.15,
  legend style={at={(0.98,0.98)},anchor=north east}]
\addplot[name path=hi,draw=none,forget plot] table[x=n_stations,y=relative_rmse_q90]{\aggn};
\addplot[name path=lo,draw=none,forget plot] table[x=n_stations,y=relative_rmse_q10]{\aggn};
\addplot[band=ewfC1,forget plot] fill between[of=hi and lo];
\addplot[ewfC1,line width=1.4pt] table[x=n_stations,y=relative_rmse]{\aggn}; \addlegendentry{observed (random portfolios)}
\addplot[ewfInk,line width=1.1pt,dashed] table[x=n_stations,y=theory]{\aggn}; \addlegendentry{Proposition~\ref{prop:smooth}}
\addplot[reference,dotted,domain=1:62,samples=2,forget plot] {\valAggTheoryLimit};
\node[font=\scriptsize,text=ewfInkSecondary,anchor=south east] at (axis cs:61,\valAggTheoryLimit) {limit $\sqrt{\bar r}=\valAggTheoryLimit$};
\end{groupplot}
\end{tikzpicture}
\caption{\textbf{Local errors average out; shared errors set a floor.} Day-ahead errors of 10\,m wind speed of the bias-corrected multi-model mean, June--July 2025. (a) Correlation of errors for all \valAggNumPairs{} station pairs against distance, with the mean per 100\,km bin. (b) RMSE of the average error of $N$ randomly chosen stations relative to the mean single-station RMSE (line: mean of 300 draws; band: 10th--90th percentile), and the prediction of Proposition~\ref{prop:smooth} with the observed mean correlation $\bar r=\valAggMeanCorr$.}
\label{fig:aggregation}
\end{figure*}


Proposition~\ref{prop:smooth} predicts how much a portfolio gains from geographic spread. The day-ahead wind errors of the bias-corrected multi-model mean are correlated at \valAggCorrNear{} on average for stations less than 100\,km apart and at \valAggCorrFar{} beyond 500\,km, with a mean of \valAggMeanCorr{} over all \valAggNumPairs{} pairs (Figure~\ref{fig:aggregation}a). Averaging the errors of all \valNumStations{} stations leaves \valAggRelRmseAll{} of the single-station RMSE, a reduction of \valAggRelRmseAllPct\,\%, and the equal-variance formula~\eqref{eq:smooth} matches the empirical curve at every portfolio size (Figure~\ref{fig:aggregation}b). With so low a mean correlation the floor $\sqrt{\bar r}=\valAggTheoryLimit$ is reached only slowly. Errors of 10\,m wind at single stations are dominated by local effects; hub-height wind and irradiance errors over a region are more strongly correlated, so real portfolios reach their floor sooner.

\subsection{From wind error to power error and to the bid}

% Figure 7: power-curve amplification and bidding.
% Data: data/processed/fig07a_power_curve.csv, fig07b_partial_load_elasticity.csv, fig07c_bid_cost_vs_quantile.csv
% Script: code/06_wind_power_and_bidding.py
\begin{figure*}[t]
\centering
\readcsv{\datadir/fig07a_power_curve.csv}\pcurve
\readcsv{\datadir/fig07b_partial_load_elasticity.csv}\elast
\readcsv{\datadir/fig07c_bid_cost_vs_quantile.csv}\bidcost
\begin{tikzpicture}
\begin{groupplot}[group style={group size=3 by 1, horizontal sep=1.3cm},
  paper, width=0.345\textwidth, height=5.0cm]
% ---------------- (a) power curve
\nextgroupplot[title={(a) Power curve, \valTurbine}, xlabel={Hub-height wind speed (m/s)},
  ylabel={Power (share of rated)}, xmin=0, xmax=26, ymin=0, ymax=1.1, xtick={0,3,6,9,12,15,20,25},
  legend style={at={(0.98,0.30)},anchor=south east}]
\fill[ewfPanel] (axis cs:3,0) rectangle (axis cs:12,1.1);
\node[font=\scriptsize,text=ewfInkSecondary,align=center] at (axis cs:7.5,0.93) {partial\\load};
\node[font=\scriptsize,text=ewfInkSecondary,align=center] at (axis cs:18.5,0.86) {full load};
\addplot[ewfC1,line width=1.4pt] table[x=wind_speed_hub_ms,y=power_normalised]{\pcurve}; \addlegendentry{manufacturer curve}
\addplot[ewfInk,dashed,line width=1.0pt,domain=3:12,samples=40] {(x^3-27)/(1728-27)}; \addlegendentry{$k(v^3-v_{ci}^3)$}
% ---------------- (b) elasticity
\nextgroupplot[title={(b) Power error per unit of wind error}, xlabel={Elasticity $\eta$ at partial load},
  ylabel={Share of forecast pairs}, xmin=2.45, xmax=5.05, ymin=0, ymax=0.27, ybar, bar width=4.2pt,
  yticklabel style={/pgf/number format/.cd,fixed,precision=2}, grid=none, ymajorgrids]
\addplot[fill=ewfC1,draw=none] table[x=elasticity_bin_mid,y=share_of_pairs]{\elast};
\draw[ewfInk,dashed,line width=1.0pt] (axis cs:3,0) -- (axis cs:3,0.27);
\node[font=\scriptsize,text=ewfInk,anchor=north west,align=left] at (axis cs:3.05,0.265) {Proposition~\ref{prop:cubic}: $\eta\ge3$\\for a cubic curve};
\node[font=\scriptsize,text=ewfInkSecondary,anchor=north east,align=right] at (axis cs:5.0,0.13) {median \valElasticityMedian};
% ---------------- (c) bidding
\nextgroupplot[title={(c) Expected cost of the bid (July 2025)}, xlabel={Quantile $\tau$ bid},
  ylabel={Cost per unit of capacity}, xmin=0.05, xmax=0.95, ymin=0.06, ymax=0.32,
  yticklabel style={/pgf/number format/.cd,fixed,precision=2},
  legend style={at={(0.02,0.98)},anchor=north west,font=\tiny}]
\addplot[ewfC1,line width=1.4pt] table[x=tau,y=cost_asymmetric]{\bidcost}; \addlegendentry{shortfall cost 3, surplus 1}
\addplot[ewfC2,line width=1.4pt] table[x=tau,y=cost_symmetric]{\bidcost}; \addlegendentry{equal costs}
\addplot[ewfC1,dashed,line width=0.8pt,forget plot] table[x=tau,y=cost_asymmetric_point_forecast]{\bidcost};
\addplot[ewfC2,dashed,line width=0.8pt,forget plot] table[x=tau,y=cost_symmetric_point_forecast]{\bidcost};
\draw[ewfInkMuted,dotted,line width=0.9pt] (axis cs:\valTauStar,0.06) -- (axis cs:\valTauStar,0.215);
\draw[ewfInkMuted,dotted,line width=0.9pt] (axis cs:0.5,0.06) -- (axis cs:0.5,0.215);
\node[font=\scriptsize,text=ewfInkSecondary,anchor=south] at (axis cs:\valTauStar,0.215) {$\tau^*=\valTauStar$};
\node[font=\scriptsize,text=ewfInkSecondary,anchor=south] at (axis cs:0.5,0.215) {$0.5$};
\end{groupplot}
\end{tikzpicture}
\caption{\textbf{The power curve amplifies wind errors, and asymmetric costs call for a quantile bid.} Pseudo power from observed and forecast 10\,m winds scaled to 120\,m (bias-corrected multi-model mean, day-ahead). (a) The turbine's power curve and its cubic idealisation between cut-in and rated speed. (b) Ratio of logarithmic power error to logarithmic wind-speed error for \valElasticityN{} pairs in which forecast and observation are both at partial load; values outside 2.5--5 are counted in the end bins. (c) Expected test cost of bidding the $\tau$-quantile of conditional quantile forecasts (solid) against bidding the point forecast (dashed); the dotted lines mark the optima predicted by Proposition~\ref{prop:bid}.}
\label{fig:power}
\end{figure*}


To isolate the effect of the power curve, observed and forecast 10\,m winds were scaled to a pseudo hub height of 120\,m with a neutral power law and passed through the power curve of a \valTurbine{} turbine (\valRatedKw\,kW) from the OpenEnergy turbine library~\cite{windpowerlib}. The resulting pseudo power is not the output of any real farm, but forecast and observation are transformed identically.

\para{Amplification (Proposition~\ref{prop:cubic}).} Where both forecast and observation lie in the partial-load range, the ratio of logarithmic power error to logarithmic wind error has a median of \valElasticityMedian{} (interquartile range \valElasticityQOne--\valElasticityQThree, \valElasticityN{} pairs), and \valElasticityShareAboveThree\,\% of pairs reach the bound of 3 (Figure~\ref{fig:power}b); the remainder lie just below rated speed, where the real power curve bends over and the cubic idealisation no longer holds. Over all conditions, including calm hours below cut-in, the relative MAE of power is \valAmplificationDA{} times that of wind speed (\valRelPowerErrDA\,\% vs.\ \valRelWindErrDA\,\% at day-ahead), and the day-ahead MAE is \valPowerMaeDA\,\% of capacity at a mean capacity factor of \valCapFactorObs\,\%.

\para{Bidding (Proposition~\ref{prop:bid}).} Day-ahead quantile forecasts of pseudo power were built from conditional empirical quantiles: for each class of forecast power, the June runs (\valBidTrainN{} cases) give the quantiles of the observed power, which are then bid on the July runs (\valBidTestN{} cases). With a shortfall three times as costly as a surplus, Proposition~\ref{prop:bid} predicts the optimal quantile $\tau^*=c_l/(c_l+c_s)=\valTauStar$. The test cost is lowest near $\tau=\valTauBest$, the cost at $\tau^*$ is within \valCostStarAboveBestPct\,\% of that minimum, and bidding the $\tau^*$ quantile instead of the point forecast lowers the expected cost by \valCostSavingPct\,\% (Figure~\ref{fig:power}c). With symmetric costs the median and the point forecast perform almost identically, as they should. The asymmetry of real imbalance prices is therefore a direct, measurable reason to produce probabilistic forecasts.
